
\def \Subject {تفسیرپذیری}
\def \Session { یک}
% \setcounter{chapter}{\Session-1}
\renewcommand{\chaptername}{سوال}
\chapter{\Subject}
\section{سوالات توضیحی}
\subsection{ \lr{PDP}}

\begin{enumerate}
\item 
در 
\lr{Partial Dependence}
در حال بررسی تاثیر ویژگی‌های 
\lr{S}
روی خروجی مدل هستیم. و میانگین خروجی مدل برای مقادیر مختلف ویژگی‌های 
\lr{S}
بدست می‌آید. حال اگر ویژگی‌های 
\lr{S}
با 
\lr{C}
رابطه داشته باشند، برای مثال در حالت افراطی وقتی که یک ویژگی دوبار تکرار شده باشد و یکی در 
\lr{S}
و یکی در
\lr{C}
باشد،  در حالات زیادی ورودی مدل دو مقدار برای آن ویژگی است که در دنیای واقعی رخ نمی‌دهد و باعث می‌شود مدل مجبور به برون‌یابی
\LTRfootnote{Extrapolate}
 خروجی معنادار نباشد. 
 
 برای مثال اگر در یک اپلیکیشن درخواست سفر، برای پیش‌بینی قیمت سفر، یک ویژگی مدت زمان سفر باشد و ویژگی دیگر فاصله بین مبدا و مقصد باشد،‌مدل برای  یافتن حالتی که مبدا و مقصد نزدیک هم هستند ولی زمان سفر خیلی زیاد است مجبور است برون‌یابی کند و احتمالا خروجی بی‌معنی بدهد که برای هوش‌مصنوعی توضیح پذیر این موضوع خوب نیست.
 

\item 
در روش 
\lr{ALE}
مقادیر مختلف ویژگی به بازه‌های کوچکی شکسته می‌شوند و سپس برای هر کدام از این بازه‌های کوچک، مقدار بیشینه و کمینه جایگزین می‌شود (بدون تغییر سایر ویژگی‌ها) و تاثیر ویژگی در بازه حساب می‌شود. سپس از مقدار کم به زیاد مقادیر جمع زده می‌شود تا تاثیر تجمعی حساب شود. سپس میانگین این مقادیر را صفر می‌کنیم و نمودار آن رسم می‌شود.

\item 
خیر. با توجه به اینکه مدل 
\lr{PDP}
اصلا به 
هم‌بستگی 
\LTRfootnote{Correlation}
اهمیت نمی‌دهد و تاثیر بود و نبود ویژگی را در نظر می‌گیرد، پس اگر مدل نسبت به یک ویژگی 
\lr{robust}
باشد، به دلیل افزونگی 
\LTRfootnote{Redundancy}
ممکن است با تغییر یک ویژگی خیلی تغییری رخ ندهد ولی اگر چند ویژگی باهم تغییر کنند تغییر کند.
\end{enumerate}

\subsection{ \lr{LIME}}
اگر مدل مورد بررسی یک مدل مقاوم
\LTRfootnote{Robust}
نباشد و مرز تصمیم آن بسیار به داده‌ها نزدیک و آشفته باشد، بسته به اینکه موقع نمونه‌برداری از فضای اطراف داده، چه نمونه‌هایی انتخاب شوند، ممکن است مدل خطی دو خط کاملا متفاوت را بیابد و ضرایب گویا نباشند. اما اگر مدل یک مدل مقاوم 
\lr{Lipchitz}
با ضریب 
\lr{L}
باشد، آنگاه از آنجایی که مدل تغییرات ناگهانی ندارد، پس مرز تصمیم مقدار پایدارتری دارد و با یک نمونه‌برداری دور نمونه (با توجه به ضریب 
\lr{L})
می‌توان دید خوبی از درون مدل یافت.

\subsection{\lr{Shapely}}
\begin{enumerate}
\item 
در محاسبه مقدار 
\lr{shapely value}
خروجی‌های مختلف مدل را به ازای بود و نبود ویژگی مد نظر حساب می‌کنیم و میانگین همه، مقدار
\lr{shapely value}
را نشان می‌دهند. اگر به فرمول داده‌شده نگاه کنیم و مقدار 
$f(x^i)$
را درنظر بگیریم، با توجه به ویژگی‌ها در فرمول مدل خطی، روی هم تاثیر نمی‌گذارند، می‌توان از تاثیر سایر ویژگی‌ها صرفنظر کرد و فقط دید عبارت مربوط به ویژگی 
\lr{j}
ام چه تغییری می‌کند که با استفاده از فرمول داده شده داریم:
\begin{latin}
	~\\
	$
	f(x^i) - E(f(x)) = \beta_j x^i_j - \beta_j E(X_j) = \beta_j (x^i_j - E(x_j)) 
	$
\end{latin}
\item 
اگر فرض کنیم در یک مدل ویژگی‌ها از هم مستقل هستند، بجای بررسی 
$2^n$
حالت مختلف، می‌توان صرفا با بررسی دو حالت به تاثیر ویژگی مدنظر رسید و مقدار 
\lr{Shapely value}
آن را بدست آورد. 

درواقع اگر ویژگی‌ها مستقل باشند، دیگر نیاز نیست احتمال شرطی‌ها حساب شود و اگر برای دو حالت بود و نبود ویژگی در حالتی که سایر ویژگی‌ها همه باشند تاثیر حساب شود، چون استقلال داریم می‌توان گفت:

$
P(A | B) = P(A)
$
پس یعنی احتمالات شرطی نیز برابر با همین مقدار می‌شود و نیازی به حساب آن‌ها و سپس میانگین گرفتن نیست.
\end{enumerate}


\subsection{\lr{NAM}}
این معماری پیشنهاد می‌کند که هر ویژگی بصورت مجزا به یک شبکه عصبی داده شود و بعد خروجی این شبکه‌ها باهم (همراه با 
\lr{bias}
)
 جمع زده شوند و خروجی را بدهند.  به اینصورت می‌توان مدلی قدرتمند‌تر از برای مثال
 \lr{logistic regression}
 داشت زیرا در 
 \lr{logistic regression}
 هر ویژگی تنها یک ضریب می‌گیرد؛ ولی در 
 \lr{NAM}
 هر ویژگی از یک شبکه عصبی عمیق رد می‌شود و شبکه می‌تواند ویژگی‌های پیچیده‌تری بیابد. 
 
 همانطور که در مقاله ذکر شده است،‌
 \lr{NAM}
 زیرمجموعه یک خانواده بزرگ‌تر از مدل‌ها به اسم 
 \lr{Generalized Additive Model (GAM)}
 است. یکی از فرض‌های اصلی در این مدل‌ها، فرض مستقل بودن ویژگی‌ها است. پس اگر ویژگی‌ها مستقل نباشند و برهم‌کنش این ویژگی‌ها خروجی را تعیین کند، 
 \lr{NAM}
 نمی‌تواند دقت خوبی بدهد. 
 
 همچنین هرچقدر تعداد ویژگی‌ها بیشتر شود، 
 \lr{NAM}
 تفسیر پیچیده‌تری از ویژگی می‌دهد که می‌تواند باعث مشکل شود. یک راه‌حلی که در مقاله ارائه می‌شود، استفاده از ویژگی‌های پیچیده‌تر است (برای مثال بجای عکس، لبه‌های آن داده شود یا از یک شبکه پیچشی گذر کند و سپس ورودی داده شود.) که در اینصورت دقت مدل بهبود می‌یابد و نیاز به پیچیدگی خود مدل 
 \lr{NAM}
 کمتر می‌شود ولی به شرط اینکه ویژگی‌های ورودی آن خود تفسیرپذیر باشند. 
 
 در \lr{NAM}
چون فضای توزیع توصیف می‌شود (و نه فقط فضای نمونه‌ای) همچنان می‌توان دید در مواردی مدل مجبور به برون‌یابی
\LTRfootnote{extrapolate}
می‌شود. ولی وزن این نواحی در نمودارها با شدت رنگ نشان داده شده‌اند و می‌توان در نمودار دید که کجا داده کمتر بوده است پس این مسئله نسبت به 
\lr{PDP}
که اصلا توجهی به توزیع داده ندارد بهبود دارد.

\section{سوالات پیاده‌سازی}
\subsection{مقداردهی}
\begin{figure}[H]
	
	\centering
	\includegraphics[width=\textwidth]{images/q1_cat.png}
	\caption{
		عکس گربه مورد بررسی
	}
	\label{q1:cat}
\end{figure}
مدل با احتمال 
$13.68$
درصد، عکس را متعلق به دسته ۲۸۲ می‌داند. 

\subsection{بخش‌بندی تصویر}
کد داده‌شده با استفاده از الگوریتم
\lr{K means}
در هر مرحله، ابتدا با محاسبه فاصله بین نزدیک‌ترین مرکز دسته و هر پیکسل‌ها تعیین می‌کند هر پیکسل متعلق به چه دسته‌ای است، میانگین مختصات اعضای هردسته را می‌یابد و مرکز دسته را به آنجا تغییر می‌دهد.  عملا این‌کار یک 
\lr{color segmentation}
است.

\begin{figure}[H]
	
	\centering
	\includegraphics[width=\textwidth]{images/q1_segcat.png}
	\caption{
		عکس گربه که به ۱۸ بخش شکسته شده است.
	}
	\label{q1:cat_seg}
\end{figure}
\subsection{تولید نمونه‌های تغییر یافته}
\begin{figure}[H]
	
	\centering
	\includegraphics[width=\textwidth]{images/q1_maskcat.png}
	\caption{
		چند نمونه از تصاویر تغییریافته.
	}
	\label{q1:cat_mask}
\end{figure}

\subsection{وزن‌دهی}
(برای فهم بهتر سوال از هوش مصنوعی کمک گرفته شد.)



	این وزن‌دهی به این دلیل رخ می‌دهد که هرچقدر از تصویر اصلی دورتر می‌شویم، تصاویر وزن کم‌تری بگیرند. به این صورت هنگام آموزش مدل خطی نیز تاثیر کمتری می‌گذارند. زیرا اگر با تغییر کمی در تصویر خروجی مدل خیلی تغییر کند، پس احتمالا ویژگی‌هایی که تغییر کردند اهمیت بیشتری برای مدل داشتند. 
\subsection{آموزش مدل خطی}
با آموزش مدل خطی، در تصویر 

می‌توان دید که مدل به اکثر نواحی توجه کمی دارد ولی نواحی که بدن و سر گربه را تشخیص می‌دهند وزن بیشتری می‌گیرند.
\begin{figure}[H]
	
	\centering
	\includegraphics[width=\textwidth]{images/q1_histcat.png}
	\caption{
	هیستوگرام احتمالات گروه گربه 
	}
	\label{q1:cat_hist}
\end{figure}
\begin{figure}[H]
	
	\centering
	\includegraphics[width=\textwidth]{images/q1_histlasso.png}
	\caption{
		هیستوگرام ضرایب مدل خطی. در اکثر موارد ضریب صفر و یا نزدیک به صفر است.
	}
	\label{q1:cat_lcoef}
\end{figure}
\begin{figure}[H]
	
	\centering
	\includegraphics[width=\textwidth]{images/q1_catcoefmap.png}
	\caption{
		خروجی 
		\lr{heatmap}
		مقدار ضرایب مدل خطی.
	}
	\label{q1:cat_lime}
\end{figure}
\subsection{تصویر خرس قطبی}
در 
\lr{heatmap}
تصویر مربوط به خرس قطبی، خروجی به اندازه تصویر گربه تمیز نیست. مدل تمرکز بیشتری روی سر و بخشی از بدن خرس دارد ولی روی پس‌زمینه و سایر موارد نیز تمرکز دارد. این موضوع خوب نیست و احتمالا به سوگیری مدل اشاره دارد (برای مثال در پس‌زمینه برفی و درختان شمالی احتمال خرس بیشتر است.)
\begin{figure}[H]
	
	\centering
	\includegraphics[width=\textwidth]{images/q1_bearcoefmap.png}
	\caption{
		خروجی 
		\lr{heatmap}
		مقدار ضرایب مدل خطی.
	}
	\label{q1:bear_lime}
\end{figure}